\section{A Frozen Learned-Policy Closed-Loop Check}
\label{sec:learned-closed-loop}

\begin{sloppypar}
Section~\ref{sec:closed-loop} evaluates non-learned policies against the
Tier-1 production closed-loop protocol. This section reports a separate,
smaller confirmatory experiment: freezing one model trained on this
benchmark's supervision and evaluating it, in closed loop, against matched
baselines under its own held-out test protocol. \textbf{This experiment uses
an independently constructed held-out test split (the frozen model's own
training-pipeline split, applied identically across all five families,
including metacdn) and is not run on the Tier-1 protocol; its miss-ratio
values are therefore not numerically comparable, cell for cell, to
Table~\ref{tab:tier1-closed-loop}.} What is comparable, and is the point of
this section, is the ordering among policies evaluated together under one
matched, held-out protocol.
\end{sloppypar}

\textbf{Setup.} A single HistGradientBoostingRegressor was trained to predict
the canonical $H=16$ \texttt{y\_loss} target, selected by an offline
validation-set regret/MAE criterion before any closed-loop replay was run,
and frozen (no retraining, retuning, feature change, or post-hoc adjustment
at any later point). It was evaluated by closed-loop replay -- evict the
resident candidate with the lowest predicted \texttt{y\_loss} at every
full-cache miss -- on the same 10 predetermined family$\times$capacity cells
(5 families $\times$ capacities 32, 128) used elsewhere in this paper, each
scored on a fixed, pre-registered held-out test interval, against matched
LRU, MRU, SIEVE, LFU (each one deterministic run) and random (20 seeds)
baselines under identical warmup (empty cache and empty bounded-history
state at the test boundary) and scoring conventions.

\begin{table}[H]
  \centering
  \footnotesize
  \caption{Frozen learned-policy closed-loop check: win/tie/loss count for the learned policy against each comparator across the same 10 predetermined cells (5 families $\times$ capacities 32, 128), under a held-out protocol distinct from Table~\ref{tab:tier1-closed-loop}. A tie is exact floating-point equality; every observed tie is one of the two wiki2018 cells, where all 6 policies score miss ratio 1.0. Macro miss ratio is the unweighted mean over the 10 cells (lower is better).}
  \label{tab:learned-closed-loop}
  \begin{tabularx}{\columnwidth}{@{}lrrrr@{}}
    \toprule
    comparator & learned wins & ties & learned losses & macro miss ratio (comparator) \\
    \midrule
    LRU    & 1 & 2 & 7 & 0.7117 \\
    MRU    & 8 & 2 & 0 & 0.8484 \\
    SIEVE  & 1 & 2 & 7 & 0.7096 \\
    LFU    & 5 & 2 & 3 & 0.7825 \\
    random (20-seed mean) & 1 & 2 & 7 & 0.7228 \\
    \midrule
    \multicolumn{4}{l}{learned macro miss ratio} & 0.7629 \\
    \bottomrule
  \end{tabularx}
\end{table}


\textbf{Result.} Table~\ref{tab:learned-closed-loop} reports the outcome the
frozen policy did \emph{not} achieve: outside the two fully degenerate
wiki2018 cells (miss ratio 1.0 for every policy, as in
Table~\ref{tab:tier1-closed-loop}), the learned policy beat LRU, SIEVE, and
the random-seed mean in exactly one of the 10 cells each -- the same cell,
metakv/cap128, in all three comparisons, by a margin of 0.0009--0.0019
absolute miss ratio -- and lost in the other 7. It beat MRU in every
non-degenerate cell (8/8) and beat LFU more often than not (5 wins, 3
losses). Excluding wiki2018 widens rather than narrows the learned-vs-LRU/
SIEVE/random gap (macro miss ratio 0.7036 for learned vs.\ 0.6397 LRU,
0.6371 SIEVE, 0.6536 random-mean on the remaining 8 cells), so the
comparison is not an artifact of averaging in the degenerate cells. In the
7 loss cells the gap is often large and not confined to the random-seed
mean: at twemcache/cap128, learned (miss ratio 0.3909) loses to every one
of the 20 random seeds (range 0.1586--0.1709), not just their mean.

\textbf{Reading this result.} This shows that, for this one frozen model
under this held-out protocol, offline predictive structure on the
finite-horizon counterfactual target did not translate into consistent
closed-loop improvement over strong non-learned baselines. It does not, by
itself, identify which mechanism is responsible: the offline selection
criterion (validation-set regret/MAE on \texttt{y\_loss}) is evaluated
per-decision against a fixed history, while the closed-loop outcome is
path-dependent on every eviction the policy itself previously made, so
predicting the label well is a different task from inducing, through a
sequence of choices, a cache-state trajectory with fewer total misses.
Tie-degenerate supervision, action-underdetermined labels, continuation-policy
mismatch, and this predict-then-act distribution shift are all consistent
with the observed gap; none is established as the cause by this evidence
alone, and this paper does not attribute one.
